% Solucionario — Tema 2: haz gaussiano, absorción y dispersión, pulsos (problemas 8 a 17).
% Convención de Saleh en todo el tema: e^{j omega t}, onda en +z ~ e^{-jkz}, chi'' < 0 = absorción,
% F{f}(nu) = int f(t) e^{-j 2 pi nu t} dt.

\textbf{Convención del tema.} Se usa la de Saleh y Teich: dependencia temporal $e^{j\omega t}$, onda que
viaja en $+z$ como $e^{-jkz}$, $k=\beta-j\alpha/2$, susceptibilidad $\chi=\chi'+j\chi''$ con
$\chi''<0$ en un medio absorbente, transformada de Fourier $A(f)=\int A(t)e^{-j2\pi ft}\dd t$
(la misma de \texttt{numpy.fft.fft}). Para los pulsos, $\tau$ es la constante de $1/e$ de la
\emph{amplitud}: $\tau_{\rm FWHM}=\sqrt{2\ln2}\,\tau=1.177\,\tau$ (Saleh, ec.~22.1-8, p.~942).
Constantes: $c_0=\SI{2.998e8}{m/s}$, $\epsilon_0=\SI{8.854e-12}{F/m}$, $e=\SI{1.602e-19}{C}$,
$m_e=\SI{9.109e-31}{kg}$.

%=====================================================================
\sol{8}{Parámetros de un haz}{Saleh Prob.~3.1-6, p.~100; §3.1B, ecs.~3.1-11 a 3.1-22, pp.~77--80}\label{sol:8}

\emph{Datos:} $\lambda=\SI{1.06}{\um}$, $P=\SI{1}{W}$, divergencia \emph{total} $2\theta_0=\SI{1}{mrad}$,
luego el semiángulo es $\theta_0=\SI{0.5}{mrad}$.

\inc{a} \textbf{Radio de la cintura.} La divergencia de campo lejano es (Saleh, ec.~3.1-20, p.~80)
\[
\theta_0=\frac{W_0}{z_0}=\frac{\lambda}{\pi W_0}
\quad\Longrightarrow\quad
W_0=\frac{\lambda}{\pi\theta_0}=\frac{\SI{1.06e-6}{m}}{\pi\,(\SI{5e-4}{})}=\SI{6.748e-4}{m}.
\]
\resultado{$W_0=\SI{0.675}{mm}$ (diámetro de mancha $2W_0=\SI{1.35}{mm}$).}

\inc{b} \textbf{Profundidad de foco.} Por (3.1-11), $z_0=\pi W_0^2/\lambda$, o más rápido
$z_0=W_0/\theta_0=\SI{6.748e-4}{m}/\SI{5e-4}{}=\SI{1.350}{m}$. La profundidad de foco es
$2z_0$ (ec.~3.1-22, p.~80):
\resultado{$2z_0=2\pi W_0^2/\lambda=\SI{2.70}{m}$.}
Orden de magnitud: un haz milimétrico en el IR cercano permanece colimado unos metros, como se espera
($z_0\sim W_0^2/\lambda\sim\SI{e-6}{m^2}/\SI{e-6}{m}$).

\inc{c} \textbf{Intensidad máxima.} La potencia es $P=\tfrac12I_0\pi W_0^2$ (ec.~3.1-15, p.~78),
luego
\[
I_0=\frac{2P}{\pi W_0^2}=\frac{2\,(\SI{1}{W})}{\pi\,(\SI{6.748e-4}{m})^2}=\frac{2}{\SI{1.4305e-6}{m^2}}
=\SI{1.40e6}{W/m^2}.
\]
\resultado{$I_0=\SI{1.40e6}{W/m^2}=\SI{140}{W/cm^2}$ (en el centro de la cintura).}

\inc{d} \textbf{Intensidad en el eje a $z=\SI{1}{m}$.} Por (3.1-13), p.~78,
\[
I(0,z)=\frac{I_0}{1+(z/z_0)^2}=\frac{\SI{140}{W/cm^2}}{1+(1/1.350)^2}=\frac{140}{1.549}\ \si{W/cm^2}.
\]
\resultado{$I(0,\SI{1}{m})=\SI{90.3}{W/cm^2}$, es decir el 64\,\% del pico, porque $z<z_0$.}

\idea{Cuidado con ``divergencia'': si el enunciado da el ángulo \emph{total} $2\theta_0$, divida entre
2 antes de usar $\theta_0=\lambda/\pi W_0$; equivocarse aquí da $W_0$ al doble y $I_0$ a la cuarta parte.
El factor $2$ en $I_0=2P/\pi W_0^2$ viene de $\int_0^\infty e^{-2\rho^2/W^2}2\pi\rho\dd\rho=\pi W^2/2$.
En el eje la intensidad cae a la mitad justo en $z=\pm z_0$.}

%=====================================================================
\sol{9}{Identificación de un haz con dos anchos}{Saleh Prob.~3.1-7, p.~100; ecs.~3.1-8 y 3.1-11, p.~77}\label{sol:9}

\emph{Planteamiento.} Sea $z_1$ la posición del primer plano medida desde la cintura (positiva si
el plano está después de la cintura) y $z_2=z_1+d$ la del segundo. Con (3.1-8) y (3.1-11),
\[
W^2(z)=W_0^2\left[1+\frac{z^2}{z_0^2}\right]=W_0^2+\frac{b\,z^2}{W_0^2},
\qquad b\equiv\left(\frac{\lambda}{\pi}\right)^2,
\]
porque $z_0=\pi W_0^2/\lambda$. Llamando $a\equiv W_0^2$, las dos medidas dan
\[
W_1^2=a+\frac{b\,z_1^2}{a},\qquad W_2^2=a+\frac{b\,(z_1+d)^2}{a}.
\]
\textbf{Paso 1: eliminar $z_1$.} Restando,
$W_2^2-W_1^2=\frac{b}{a}\big[(z_1+d)^2-z_1^2\big]=\frac{b}{a}(2z_1d+d^2)$, de donde
\[
z_1=\frac{1}{2d}\left[\frac{a\,\Delta}{b}-d^2\right],\qquad \Delta\equiv W_2^2-W_1^2 .
\]
\textbf{Paso 2: ecuación para $a$.} Multiplicando la primera ecuación por $a$:
$a^2-aW_1^2+bz_1^2=0$. Sustituyendo $z_1$,
\[
bz_1^2=\frac{b}{4d^2}\left[\frac{a^2\Delta^2}{b^2}-\frac{2a\Delta d^2}{b}+d^4\right]
=\frac{a^2\Delta^2}{4bd^2}-\frac{a\Delta}{2}+\frac{bd^2}{4},
\]
y queda una \emph{cuadrática}:
\[
\underbrace{\left(1+\frac{\Delta^2}{4bd^2}\right)}_{A}a^2-\underbrace{\left(W_1^2+\frac{\Delta}{2}\right)}_{B}a+\underbrace{\frac{bd^2}{4}}_{C}=0 .
\]
\textbf{Paso 3: números.} $\lambda=\SI{10.6}{\um}$: $b=\SI{1.1385e-11}{m^2}$;
$W_1^2=\SI{2.8866e-6}{m^2}$, $W_2^2=\SI{1.14244e-5}{m^2}$, $\Delta=\SI{8.5378e-6}{m^2}$, $d=\SI{0.1}{m}$:
\[
A=161.07,\qquad B=\SI{7.1555e-6}{m^2},\qquad C=\SI{2.846e-14}{m^4},
\]
\[
a=\frac{B\pm\sqrt{B^2-4AC}}{2A}=\frac{(7.1555\pm5.7326)\times10^{-6}}{322.15}\ \si{m^2}.
\]
\begin{itemize}
\item Raíz $+$: $a=\SI{4.001e-8}{m^2}$, $W_0=\SI{0.200}{mm}$, $z_0=\pi a/\lambda=\SI{1.19}{cm}$ y
\[z_1=\frac{1}{0.2\,\mathrm{m}}\left[\frac{(4.001\times10^{-8})(8.538\times10^{-6})}{1.1385\times10^{-11}}-0.01\right]\mathrm{m^2}=\SI{+10.0}{cm}.\]
Los dos planos ($z_1=10$ cm, $z_2=20$ cm) están del mismo lado, después de la cintura.
\item Raíz $-$: $a=\SI{4.42e-9}{m^2}$, $W_0=\SI{66.5}{\um}$, $z_0=\SI{1.31}{mm}$, $z_1=\SI{-3.34}{cm}$,
$z_2=\SI{+6.66}{cm}$: la cintura queda \emph{entre} los dos planos.
\end{itemize}
Se comprueba que ambas reproducen exactamente $W_1=\SI{1.699}{mm}$ y $W_2=\SI{3.380}{mm}$.

\resultado{Solución con los dos planos del mismo lado: cintura \SI{10.0}{cm} antes del primer plano,
$W_0=\SI{0.200}{mm}$, $z_0=\SI{1.19}{cm}$. Existe una segunda solución matemáticamente válida,
$W_0=\SI{66.5}{\um}$ con la cintura \SI{3.34}{cm} después del primer plano (entre ambos).}

\textbf{[N] Solución numérica y unicidad.} El programa resuelve la cuadrática y además lanza
\texttt{fsolve} desde 36 puntos de partida; solo encuentra esas dos soluciones.
\codigo{t2_p09_dos_anchos.py}
Salida relevante:\\ \texttt{W0 = 0.2000 mm, z0 = 1.1857 cm, z1 = +10.002 cm} y\\
\texttt{W0 = 0.0665 mm, z0 = 0.1309 cm, z1 = -3.344 cm}.
\begin{figure}[H]\centering\includegraphics[width=0.75\linewidth]{t2_p09_dos_anchos.png}
\caption{Problema 9: las dos soluciones $W(z)$ que pasan por los dos anchos medidos (puntos).}\end{figure}
\emph{Discusión.} Dos anchos no fijan el haz de forma única, porque $W(z)$ es par en $z$ respecto
a la cintura: un haz muy enfocado ($z_0$ pequeño) cuya cintura está entre los planos también pasa por
los dos puntos. Para decidir hace falta información adicional: el radio de curvatura en un plano
(que da el signo de $z$, Ejercicio~3.1-3, p.~84), un tercer ancho, o saber que la cintura no está entre
los planos. Esta última es la suposición implícita del libro, y con ella la solución es única
($W_0=\SI{0.2}{mm}$), como dice la respuesta publicada.

\idea{El problema inverso del haz gaussiano se reduce a una cuadrática en $W_0^2$; siempre revise
las dos raíces. Datos $(W,R)$ en un plano sí dan solución única (3.1-25)--(3.1-26); dos anchos no.}

%=====================================================================
\sol{10}{Enfoque y refracción de un haz}{Saleh Probs.~3.2-6, 3.2-7, 3.2-8, pp.~100--101; §3.2A--B, ecs.~3.2-5 a 3.2-15, pp.~87--89}\label{sol:10}

\inc{a} \textbf{Focal mínima.} Datos: $\lambda=\SI{488}{nm}$, $W_0=\SI{0.5}{mm}$, mancha de diámetro
\SI{100}{\um}, es decir $W_0'=\SI{50}{\um}$. El rango de Rayleigh incidente es
\[
z_0=\frac{\pi W_0^2}{\lambda}=\frac{\pi(\SI{5e-4}{m})^2}{\SI{4.88e-7}{m}}=\SI{1.609}{m}.
\]
\emph{Con la ley del parámetro $q$.} En la cintura $q_1=jz_0$ (ec.~3.1-5). Si la lente está a
distancia $z$ de la cintura, justo antes de ella $q=z+jz_0$, y tras la lente
$1/q_2=1/q-1/f$. Escribiendo $q_2=z'_{\rm rel}+jz_0'$ y separando partes real e imaginaria se llega a
la transformación de parámetros (3.2-5)--(3.2-9), p.~87: $W_0'=MW_0$ con
\[
M=\frac{M_r}{\sqrt{1+r^2}}=\frac{|f/(z-f)|}{\sqrt{1+z_0^2/(z-f)^2}}=\frac{|f|}{\sqrt{(z-f)^2+z_0^2}} .
\]
Para $f$ fijo, $M$ es máximo cuando $z=f$ (la lente a una distancia focal de la cintura), y vale
$M_{\rm max}=f/z_0$. Como se necesita reducir el haz en un factor $M=W_0'/W_0=0.1$, y $M\le f/z_0$,
\[
\frac{f}{z_0}\ \ge\ \frac{W_0'}{W_0}\quad\Longrightarrow\quad
f_{\min}=\frac{W_0'}{W_0}z_0=0.1\times\SI{1.609}{m}=\SI{16.1}{cm}.
\]
Con $f<\SI{16.1}{cm}$ la mancha \emph{siempre} sale menor que \SI{100}{\um}, sin importar dónde se ponga
la lente (en realidad lo que se busca es la focal más corta que todavía puede dar exactamente
\SI{100}{\um}; con focales más largas hay dos posiciones de la lente que lo logran). Comprobación con el haz colimado (3.2-15):
$W_0'\approx\theta_0f=\frac{\lambda}{\pi W_0}f=(\SI{3.107e-4}{})(\SI{0.161}{m})=\SI{50}{\um}$.
\resultado{$z_0=\SI{1.61}{m}$; $f_{\min}=W_0'z_0/W_0=\SI{16.1}{cm}$, con la lente a $z=f$ de la cintura
(la nueva cintura queda entonces en $z'=f$ detrás de la lente, por (3.2-6)).}

\inc{b} \textbf{Refracción en una interfaz plana.} La cintura está sobre la interfaz, donde el frente
de onda es plano ($R=\infty$). La condición de frontera exige continuidad del campo transversal, así
que el perfil $e^{-\rho^2/W_0^2}$ es el mismo a ambos lados: \emph{$W_0$ no cambia} y el frente sigue
plano, de modo que el haz transmitido también tiene su cintura en la interfaz. Lo que cambia es la
longitud de onda en el medio, $\lambda=\lambda_0/n$ (Saleh p.~77 y pista del Prob.~3.2-8). Entonces
\[
\theta_{\rm vidrio}=\frac{\lambda_0/n}{\pi W_0}=\frac{\theta_{\rm aire}}{n}=\frac{\SI{1}{mrad}}{1.5},
\qquad
z_{0,\rm vidrio}=\frac{\pi W_0^2}{\lambda_0/n}=n\,z_{0,\rm aire}.
\]
\resultado{Divergencia en el vidrio $=\SI{0.667}{mrad}$; el rango de Rayleigh crece en un factor
$n=1.5$. (Coincide con la ley de Snell para ángulos pequeños, $\theta_{\rm aire}=n\theta_{\rm vidrio}$.)}

\inc{c} \textbf{[N] Cintura enfocada en función de la distancia a la lente.} Con $z_0=\SI{50}{cm}$,
$\lambda=\SI{488}{nm}$: $W_0=\sqrt{\lambda z_0/\pi}=\SI{0.2787}{mm}$, y $f=\SI{5}{cm}$. Se grafican
$W_0'=MW_0$ (3.2-5) y $z'=f+M^2(z-f)$ (3.2-6).
\codigo{t2_p10_enfoque.py}
\begin{figure}[H]\centering\includegraphics[width=0.8\linewidth]{t2_p10_enfoque.png}
\caption{Problema 10(c): radio de la cintura enfocada $W_0'$ y distancia lente--cintura $z'$ contra la
distancia $z$ entre la cintura incidente y la lente ($z_0=\SI{50}{cm}$, $f=\SI{5}{cm}$).}\end{figure}
Números que imprime el programa:
\begin{center}\small
\begin{tabular}{cccc}\toprule
$z$ (m) & $M$ & $W_0'$ (\si{\um}) & $z'$ (cm)\\\midrule
0 & 0.0995 & 27.73 & 4.950\\
0.05 ($=f$) & 0.1000 & 27.87 (máximo) & 5.000\\
0.5 & 0.0743 & 20.72 & 5.249\\
1.0 & 0.0466 & 12.98 & 5.206\\
3.0 & 0.0167 & 4.66 & 5.082\\\bottomrule
\end{tabular}\end{center}
\emph{Lectura.} (i) Para $z\ll z_0$ la cintura enfocada es prácticamente constante,
$W_0'=W_0/\sqrt{1+(z_0/f)^2}=\SI{27.7}{\um}$ (3.2-13): el haz está colimado y la lente ``no sabe'' a qué
distancia está la cintura. (ii) Para $z-f\gg z_0$ se recupera la óptica de rayos, $W_0'\approx M_rW_0$ y
$1/z+1/z'=1/f$ (3.2-10)--(3.2-12): la cintura se comporta como un objeto que se forma en imagen. (iii) La
nueva cintura siempre está cerca del plano focal ($4.95<z'<5.25$ cm), con el máximo de $z'$ cerca de
$z=z_0+f$.

\idea{$M=f/\sqrt{(z-f)^2+z_0^2}\le f/z_0$: el aumento de un haz gaussiano nunca supera el de rayos y
es máximo con la lente a una focal de la cintura. Al pasar a un medio de índice $n$ se usa
$\lambda_0/n$ en todas las fórmulas: $W_0$ igual, $z_0\to nz_0$, $\theta_0\to\theta_0/n$. En un
experimento no lineal la intensidad en el foco es $\propto1/W_0'^2$, por eso importa este control.}

%=====================================================================
\sol{11}{Fase de Gouy}{Saleh §3.1B, ecs.~3.1-10, 3.1-23, 3.1-24, p.~81; §3.3, ec.~3.3-10, p.~96; §10.2D, ec.~10.2-33, p.~387; Prob.~3.3-4, p.~101}\label{sol:11}

\inc{a} \textbf{[N] Gráfica.} La fase del haz gaussiano (3.1-7) es
$\varphi(\rho,z)=kz-\zeta(z)+k\rho^2/2R(z)$ (3.1-23); en el eje, $\varphi(0,z)=kz-\zeta(z)$ (3.1-24), con
$\zeta(z)=\tan^{-1}(z/z_0)$ (3.1-10). Restando la onda plana $kz$ queda $\varphi(0,z)-kz=-\zeta(z)$.
\codigo{t2_p11_gouy.py}
\begin{figure}[H]\centering\includegraphics[width=0.75\linewidth]{t2_p11_gouy.png}
\caption{Problema 11: fase de Gouy $\zeta(z)$, exceso de fase axial respecto a una onda plana
$-\zeta(z)$ y la de dos modos Hermite--Gauss, para $-5z_0<z<5z_0$.}\end{figure}
Valores: $\zeta(\pm z_0)=\pm\pi/4$; $\zeta(\pm5z_0)=\pm\SI{1.373}{rad}$, así que en la ventana graficada
el salto es $\SI{2.747}{rad}=0.874\pi$; el resto hasta $\pi$ se acumula muy lentamente fuera de ella.

\inc{b} \textbf{Desfase total $\pi$.} Como $\tan^{-1}u\to\pm\pi/2$ cuando $u\to\pm\infty$,
\[
\Delta\zeta=\zeta(+\infty)-\zeta(-\infty)=\frac{\pi}{2}-\left(-\frac{\pi}{2}\right)=\pi .
\]
\resultado{Al atravesar el foco, la fase en el eje se atrasa $\pi$ respecto a la de una onda plana
(efecto Gouy, p.~81); la mitad, $\pi/2$, se acumula entre $-z_0$ y $z_0$.}
\emph{Por qué ocurre.} El haz confinado contiene componentes de onda plana inclinadas, cuya
componente axial del vector de onda es $k_z=\sqrt{k^2-k_T^2}<k$; al estrecharse en el foco crece la
dispersión de $k_T$ (principio de incertidumbre transversal) y la fase axial avanza menos que $kz$.
Signo: con $e^{j\omega t}$ y onda $e^{-jkz}$, el término $+j\zeta$ en (3.1-7) es un \emph{retraso} del
frente de onda; en la convención de Boyd ($e^{-i\omega t}$) aparece como $e^{-i\zeta}$, pero el
fenómeno es el mismo. En óptica no lineal, el término de Gouy de un haz enfocado entra en el desajuste de
fase efectivo de la generación de armónicos (Boyd §2.10).

\inc{c} \textbf{Modos Hermite--Gauss y frecuencias de un resonador.} El modo $(l,m)$ tiene fase axial
$\varphi(0,z)=kz-(l+m+1)\zeta(z)$ (ec.~10.2-31, p.~386). En un resonador de espejos esféricos en $z_1$ y $z_2$
($d=z_2-z_1$), la condición de resonancia es que la fase en un viaje de ida y vuelta sea múltiplo
de $2\pi$:
\[
2kd-2(l+m+1)\Delta\zeta=2\pi q,\qquad \Delta\zeta=\zeta(z_2)-\zeta(z_1).
\]
Con $k=2\pi\nu/c$ y $\nu_F=c/2d$ se despeja (ec.~10.2-33, p.~387)
\[
\nu_{l,m,q}=q\,\nu_F+(l+m+1)\frac{\Delta\zeta}{\pi}\,\nu_F .
\]
\resultado{Consecuencia: los modos transversales de distinto orden $l+m$ \emph{no} son degenerados en
frecuencia; se desplazan $\frac{\Delta\zeta}{\pi}\nu_F$ por cada unidad de $l+m$. Los modos con igual
$l+m$ sí tienen la misma frecuencia.}
\emph{Ejemplo (Prob.~3.3-4).} $z_0=\SI{30}{cm}$, espejos en $z=\pm z_0$: $d=\SI{60}{cm}$,
$\nu_F=c/2d=\SI{249.8}{MHz}$, $\Delta\zeta=\pi/4-(-\pi/4)=\pi/2$. La condición ``fase entre $-z_0$ y $z_0$
múltiplo de $\pi$'' da $\nu=\nu_F[q+(l+m+1)/2]$: las frecuencias permitidas son múltiplos de
$\nu_F/2=\SI{124.9}{MHz}$ (los de $l+m$ par caen en múltiplos impares de $\nu_F/2$ y los de $l+m$ impar
en múltiplos pares). En la banda $10^{14}\pm2\times10^9$ Hz hay 32 de ellas (33 si se toma
$c=\SI{3e8}{m/s}$, con el que $10^{14}$ Hz es múltiplo exacto de \SI{125}{MHz}). Es el resonador
confocal: los modos transversales coinciden o quedan a mitad de camino entre los longitudinales.

\idea{Gouy: $\zeta(z)=\tan^{-1}(z/z_0)$, $\pi/4$ en $z_0$, $\pi$ en total; para HG$_{lm}$ multiplíquese por
$l+m+1$ (para Laguerre--Gauss, $l+2m+1$). Es la razón por la que en un láser los modos transversales
producen batidos a frecuencias distintas de $\nu_F$.}

%=====================================================================
\sol{12}{Absorción débil y fuerte; medios conductores}{Saleh §5.5A, ecs.~5.5-5, 5.5-7, 5.5-8, 5.5-11, 5.5-12, pp.~172--173; §5.5D, ecs.~5.5-33 a 5.5-40, pp.~181--183}\label{sol:12}

La relación general es (Saleh, ec.~5.5-5, p.~172)
\[
n-j\frac{\alpha}{2k_0}=\sqrt{1+\chi'+j\chi''},\qquad k_0=\frac{2\pi}{\lambda_0},
\]
tomando la raíz del cuarto cuadrante ($n>0$, $\alpha>0$ cuando $\chi''<0$). A $\lambda_0=\SI{1}{\um}$,
$k_0=\SI{6.2832e6}{m^{-1}}$.

\inc{a} \textbf{$\chi=1.25-j0.01$ (absorción débil).} \emph{Exacto:} $1+\chi=2.25-j0.01$. Con
$\sqrt{x+jy}$, $|x+jy|=\sqrt{2.25^2+0.01^2}=2.2500222$, módulo de la raíz $1.5000074$, ángulo
$\tfrac12\tan^{-1}(-0.01/2.25)=\SI{-2.2222e-3}{rad}$:
\[
n=1.5000074\cos(2.2222\times10^{-3})=1.500004,\qquad
\frac{\alpha}{2k_0}=1.5000074\sin(2.2222\times10^{-3})=\SI{3.3333e-3}{}.
\]
\emph{Aproximación débil} ($|\chi''|\ll1+\chi'$; ecs.~5.5-7, 5.5-8):
$n\approx\sqrt{1+\chi'}=1.5$ y $\alpha\approx-\frac{k_0}{n}\chi''=\frac{(6.2832\times10^6)(0.01)}{1.5}$.
En ambos casos
\[
\alpha=\SI{4.189e4}{m^{-1}}=\SI{418.9}{cm^{-1}},\qquad
\alpha_{\rm dB}=10\log_{10}(e)\,\alpha=4.343\times418.9=\SI{1819}{dB/cm}.
\]
\resultado{$n=1.5000$, $\alpha=\SI{4.19e4}{m^{-1}}=\SI{419}{cm^{-1}}=\SI{1.82e3}{dB/cm}$; la aproximación
débil coincide con el resultado exacto en 4 cifras (el error relativo es del orden de
$\delta^2$, con $\delta=\chi''/(1+\chi')=0.0044$).}

\inc{b} \textbf{$\chi=-0.5-j40$ (absorción fuerte).} \emph{Exacto:} $1+\chi=0.5-j40$,
$|1+\chi|=40.003$, módulo de la raíz $6.3248$, ángulo $\tfrac12\tan^{-1}(-40/0.5)=\tfrac12(-1.5583)=\SI{-0.77915}{rad}$:
\[
n=6.3248\cos(0.77915)=4.500,\qquad \frac{\alpha}{2k_0}=6.3248\sin(0.77915)=4.444,\qquad
\alpha=2k_0(4.444)=\SI{5.585e7}{m^{-1}}.
\]
\emph{Aproximación fuerte} ($|\chi''|\gg|1+\chi'|$; ecs.~5.5-11, 5.5-12): $\sqrt{j\chi''}=\sqrt{-j}\sqrt{-\chi''}
=\frac{1-j}{\sqrt2}\sqrt{-\chi''}$, luego
\[
n\approx\sqrt{\frac{-\chi''}{2}}=\sqrt{20}=4.472,\qquad
\alpha\approx2k_0\sqrt{\frac{-\chi''}{2}}=2(6.2832\times10^6)(4.472)=\SI{5.620e7}{m^{-1}}.
\]
\resultado{Exacto: $n=4.50$, $\alpha=\SI{5.58e7}{m^{-1}}$ ($1/\alpha=\SI{17.9}{nm}$). Fuerte: $n\approx4.47$,
$\alpha\approx\SI{5.62e7}{m^{-1}}$; error de 0.6\,\%, del orden de $|1+\chi'|/|\chi''|=1/80$.}

\inc{c} \textbf{Aluminio con Drude.} Frecuencia de plasma (ec.~5.5-40, p.~183):
\[
\omega_p=\sqrt{\frac{Ne^2}{\epsilon_0m}}=\sqrt{\frac{(1.8\times10^{29})(1.602\times10^{-19})^2}{(8.854\times10^{-12})(9.109\times10^{-31})}}
\]
\[
=\sqrt{\SI{5.728e32}{s^{-2}}}=\SI{2.393e16}{rad/s},
\qquad \lambda_p=\frac{2\pi c_0}{\omega_p}=\SI{78.7}{nm}.
\]
Permitividad efectiva: de (5.5-38) con $\epsilon=\epsilon_0$ y $\sigma_0=\epsilon_0\omega_p^2\tau$
(que sale de $\omega_p^2=\sigma_0/\epsilon_0\tau$),
\[
\frac{\epsilon_e}{\epsilon_0}=1+\frac{\omega_p^2\tau}{j\omega(1+j\omega\tau)}
=1-\frac{\omega_p^2\tau^2}{1+\omega^2\tau^2}-j\,\frac{\omega_p^2\tau}{\omega(1+\omega^2\tau^2)} ,
\]
donde se multiplicó numerador y denominador por el conjugado de $j\omega-\omega^2\tau$. A
\SI{1.55}{\um}: $\omega=2\pi c_0/\lambda=\SI{1.2153e15}{rad/s}$, $\omega\tau=9.72$,
$\omega_p\tau=191.5$, $\omega_p^2\tau^2=\SI{3.666e4}{}$:
\[
\frac{\epsilon_e}{\epsilon_0}=1-\frac{36\,660}{95.5}-j\frac{36\,660/9.72}{95.5}=-383-j39.5 .
\]
A \SI{550}{nm} ($\omega\tau=27.4$) el programa da $-47.8-j1.78$. La reflectancia a incidencia normal,
$\mathcal R=|(1-\sqrt{\epsilon_e/\epsilon_0})/(1+\sqrt{\epsilon_e/\epsilon_0})|^2$, es 0.99 en ambas longitudes
de onda, y la profundidad de penetración $1/\alpha\approx\SI{6}{nm}$.
\resultado{$\omega_p=\SI{2.39e16}{rad/s}$, $\lambda_p=\SI{78.7}{nm}$ (UV de vacío),
$\epsilon_e/\epsilon_0=-383-j39.5$ a \SI{1.55}{\um}. En el visible $\omega<\omega_p$, $\mathrm{Re}\,\epsilon_e<0$:
banda prohibida de plasma, la onda es evanescente y se refleja casi toda: el Al es buen espejo del
visible al IR (y sigue siéndolo en el UV cercano, hasta $\lambda_p$).}

\inc{d} \textbf{Silicio dopado tipo n.} Ahora el fondo es $\epsilon=11.7\epsilon_0$ y la masa es $m^*$.
Para $\omega\tau\gg1$, (5.5-38) da $\epsilon_e\approx\epsilon-\sigma_0/\omega^2\tau=\epsilon_0\left(11.7-\omega_p^2/\omega^2\right)$
con $\omega_p^2=Ne^2/\epsilon_0m^*$. Con $N=\SI{e20}{cm^{-3}}=\SI{e26}{m^{-3}}$ y $m^*=0.26m_e$:
\[
\omega_p^2=\frac{(10^{26})(1.602\times10^{-19})^2}{(8.854\times10^{-12})(0.26)(9.109\times10^{-31})}=\SI{1.224e30}{s^{-2}},
\qquad \omega_p=\SI{1.106e15}{rad/s}.
\]
$\mathrm{Re}\,\epsilon_e=0$ cuando $\omega=\omega_p/\sqrt{11.7}=\SI{3.234e14}{rad/s}$:
\resultado{$\lambda=2\pi c_0\sqrt{11.7}/\omega_p=\SI{5.8}{\um}$. Para $\lambda>\SI{5.8}{\um}$ el Si tan dopado
se vuelve ``metálico''; a \SI{1.55}{\um} la contribución de los portadores libres es pequeña pero no
nula, y su parte imaginaria es la absorción por portadores libres que limita las guías de Si dopadas
(y la que generan los portadores por absorción de dos fotones).}

\codigo{t2_p12_absorcion.py}

\idea{Reconozca el régimen antes de aproximar: débil si $|\chi''|\ll1+\chi'$ ($n$ lo da $\chi'$ y
$\alpha\propto\chi''$); fuerte si $|\chi''|\gg|1+\chi'|$ ($n\approx\alpha/2k_0$). En Saleh, absorción
$\Leftrightarrow\chi''<0$ y $\mathrm{Im}\,\epsilon_e<0$; en Boyd ($e^{-i\omega t}$) los signos de las
partes imaginarias se invierten. dB $=4.343\,\alpha$ (intensidad).}

%=====================================================================
\sol{13}{Modelo de Lorentz y Kramers--Kronig numérico}{Saleh §5.5B--C, ecs.~5.5-13 a 5.5-27, pp.~175--179; §5.6, ec.~5.6-2, p.~185 y pp.~188--189}\label{sol:13}

\inc{a} \textbf{Susceptibilidad del oscilador.} Ecuación de movimiento del electrón ligado (carga $-e$):
\[
\ddot x+\sigma\dot x+\omega_0^2x=-\frac{e}{m}\mathcal E .
\]
Con $\mathcal E=\mathrm{Re}\{Ee^{j\omega t}\}$ y $x=\mathrm{Re}\{Xe^{j\omega t}\}$, cada derivada temporal
es un factor $j\omega$:
\[
(-\omega^2+j\sigma\omega+\omega_0^2)X=-\frac{e}{m}E
\quad\Longrightarrow\quad
X=-\frac{e/m}{\omega_0^2-\omega^2+j\sigma\omega}E .
\]
La polarización es $P=N(-e)X$ (momento dipolar $-ex$ por átomo):
\[
P=\frac{Ne^2/m}{\omega_0^2-\omega^2+j\sigma\omega}E\equiv\epsilon_0\chi E
\quad\Longrightarrow\quad
\chi=\frac{Ne^2}{\epsilon_0m}\frac{1}{\omega_0^2-\omega^2+j\sigma\omega}
=\chi_0\frac{\omega_0^2}{\omega_0^2-\omega^2+j\sigma\omega},\quad \chi_0=\frac{Ne^2}{\epsilon_0m\omega_0^2}
\]
(ec.~5.5-17, p.~176). Dividiendo numerador y denominador entre $(2\pi)^2$, con $\omega=2\pi\nu$ y
$\Delta\nu\equiv\sigma/2\pi$ ($\sigma\omega/4\pi^2=\nu\Delta\nu$):
\[
\chi(\nu)=\chi_0\frac{\nu_0^2}{\nu_0^2-\nu^2+j\nu\Delta\nu}\qquad\text{(5.5-19), p.~177.}
\]
Multiplicando por el conjugado del denominador:
\[
\chi'=\chi_0\frac{\nu_0^2(\nu_0^2-\nu^2)}{(\nu_0^2-\nu^2)^2+(\nu\Delta\nu)^2},\qquad
\chi''=-\chi_0\frac{\nu_0^2\,\nu\Delta\nu}{(\nu_0^2-\nu^2)^2+(\nu\Delta\nu)^2}\qquad\text{(5.5-20), (5.5-21).}
\]
\emph{Cerca de resonancia} ($|\nu-\nu_0|\ll\nu_0$): $\nu_0^2-\nu^2=(\nu_0-\nu)(\nu_0+\nu)\approx2\nu_0(\nu_0-\nu)$
y $\nu\Delta\nu\approx\nu_0\Delta\nu$, así que
\[
\chi\approx\chi_0\frac{\nu_0^2}{2\nu_0(\nu_0-\nu)+j\nu_0\Delta\nu}=\chi_0\frac{\nu_0/2}{(\nu_0-\nu)+j\Delta\nu/2}
=\chi_0\frac{\nu_0}{2}\frac{(\nu_0-\nu)-j\Delta\nu/2}{(\nu_0-\nu)^2+(\Delta\nu/2)^2}.
\]
Entonces
\[
\chi''\approx-\chi_0\frac{\nu_0\Delta\nu}{4}\frac{1}{(\nu_0-\nu)^2+(\Delta\nu/2)^2},\qquad
\chi'\approx\chi_0\frac{\nu_0}{2}\frac{\nu_0-\nu}{(\nu_0-\nu)^2+(\Delta\nu/2)^2}.
\]
$\chi''$ es una lorentziana: vale $-\chi_0\nu_0/\Delta\nu=-\chi_0Q$ en $\nu_0$ y cae a la mitad cuando
$(\nu_0-\nu)^2=(\Delta\nu/2)^2$, es decir en $\nu=\nu_0\pm\Delta\nu/2$: su FWHM es $\Delta\nu=\sigma/2\pi$.
El cociente de ambas da
\[
\frac{\chi'}{\chi''}=\frac{(\nu_0/2)(\nu_0-\nu)}{-(\nu_0\Delta\nu/4)}=\frac{2(\nu-\nu_0)}{\Delta\nu}
\quad\Longrightarrow\quad
\chi'=2\frac{\nu-\nu_0}{\Delta\nu}\chi''\qquad\text{(5.5-24), p.~178.}
\]

\inc{b} \textbf{[N] $n(\nu)-1$ y $\alpha(\nu)$.} Con $\chi_0=10^{-4}$, $\nu_0=\SI{500}{THz}$,
$\Delta\nu=\SI{5}{THz}$ ($Q=100$) y huésped $n_0=1$, se usan las fórmulas diluidas
(5.5-26)--(5.5-27), $n\approx n_0+\chi'/2n_0$ y $\alpha\approx-(2\pi\nu/n_0c_0)\chi''$, y se comparan con
el resultado exacto de (5.5-5).
\begin{figure}[H]\centering\includegraphics[width=0.8\linewidth]{t2_p13_n_alpha.png}
\caption{Problema 13(b): índice y coeficiente de absorción de un medio de Lorentz diluido
($\chi_0=10^{-4}$, $\nu_0=\SI{500}{THz}$, $\Delta\nu=\SI{5}{THz}$).}\end{figure}
Resultados: $-\chi''(\nu_0)=\chi_0Q=0.01$, $\alpha(\nu_0)=\frac{2\pi(5\times10^{14})}{2.998\times10^8}(0.01)=\SI{1.05e5}{m^{-1}}=\SI{1.05e3}{cm^{-1}}$,
FWHM de $\alpha$ igual a \SI{4.95}{THz} $\approx\Delta\nu$; $n-1$ oscila entre $+\SI{2.51e-3}{}$ (en
\SI{497.5}{THz}) y $-\SI{2.49e-3}{}$ (en \SI{502.5}{THz}), es decir $\pm\chi_0Q/4$, en $\nu_0\mp\Delta\nu/2$.
Las curvas diluida y exacta son indistinguibles porque $|\chi|\le0.01$.

\inc{c} \textbf{[N] Kramers--Kronig numérico.}
\emph{Nota importante sobre el signo.} La ec.~(5.5-13) del libro (p.~175) está impresa como
$\chi'(\nu)=\frac{2}{\pi}\int_0^\infty\frac{s\chi''(s)}{s^2-\nu^2}\dd s$. Esa es la forma correcta para la
convención $e^{-i\omega t}$ (Boyd), en la que $\chi''>0$ en un absorbente. Con la convención del propio
Saleh ($\chi''<0$) da $-\chi'$: por ejemplo, en el límite $\nu\to0$ daría
$\chi'(0)=\frac{2}{\pi}\int\chi''(s)/s\,\dd s<0$, cuando para el oscilador $\chi'(0)=\chi_0>0$. El
programa lo comprueba: la fórmula impresa da $\chi'(495)=-\SI{4.028e-3}{}$ frente al valor analítico
$+\SI{4.028e-3}{}$. La forma consistente con el resto del capítulo es
\[
\chi'(\nu)=-\frac{2}{\pi}\,\mathrm{VP}\!\int_0^\infty\frac{s\,\chi''(s)}{s^2-\nu^2}\dd s
=\frac{2}{\pi}\,\mathrm{VP}\!\int_0^\infty\frac{s\,\chi''(s)}{\nu^2-s^2}\dd s .
\]
(Lo mismo pasa con (5.5-14): ambas relaciones impresas son un par consistente para la convención
$\chi''>0$.)

\emph{Valor principal.} El integrando diverge en $s=\nu$. Se escribe
$\frac{s\chi''(s)}{s^2-\nu^2}=\frac{1}{s-\nu}\cdot\frac{s\chi''(s)}{s+\nu}$ y en un tramo
$[\nu-40,\nu+40]$ THz se usa \texttt{quad(\ldots, weight='cauchy', wvar=nu)}, que calcula
$\mathrm{VP}\int f(s)/(s-\nu)\dd s$ (por cuadratura de Clenshaw--Curtis modificada, sin evaluar el polo). El
resto de $[0,\infty)$ se integra normalmente, marcando $\nu_0$ como punto delicado por el pico
lorentziano. Truncar se hace igual pero solo sobre $|s-\nu_0|<3\Delta\nu$.
\codigo{t2_p13_lorentz_kk.py}
\begin{figure}[H]\centering\includegraphics[width=0.8\linewidth]{t2_p13_kk.png}
\caption{Problema 13(c): $\chi'(\nu)$ analítica, la obtenida por Kramers--Kronig sobre todo el
espectro y la obtenida truncando la integral en $\nu_0\pm3\Delta\nu$; abajo, el error absoluto.}\end{figure}
\begin{center}\small
\begin{tabular}{cccc}\toprule
$\nu$ (THz) & $\chi'$ analítica & KK completa & KK truncada ($\pm3\Delta\nu$)\\\midrule
400 & $+2.776\times10^{-4}$ & $+2.776\times10^{-4}$ & $+2.490\times10^{-4}$\\
495 & $+4.028\times10^{-3}$ & $+4.028\times10^{-3}$ & $+4.046\times10^{-3}$\\
600 & $-2.271\times10^{-4}$ & $-2.271\times10^{-4}$ & $-2.038\times10^{-4}$\\
1 (casi estático) & $1.000\times10^{-4}$ & $1.000\times10^{-4}$ & $0.895\times10^{-4}$\\\bottomrule
\end{tabular}\end{center}
La integral completa reproduce $\chi'$ con error relativo $10^{-7}$: $\chi$ del oscilador es causal y
cumple Kramers--Kronig exactamente. \emph{Truncada:} (i) lejos de la resonancia $\chi'$ queda un
10\,\% corta, porque se pierden las alas de la lorentziana, que decaen solo como $1/(\nu-\nu_0)^2$
(la fracción del área de $\chi''$ dentro de $\pm3\Delta\nu$ es $\frac{2}{\pi}\tan^{-1}6=0.895$, y el valor
estático sale $0.895\chi_0$); (ii) en los bordes $\nu=\nu_0\pm3\Delta\nu$ aparece una divergencia
logarítmica espuria, porque el valor principal necesita que el integrando sea suave a los dos lados
del polo; (iii) cerca de $\nu_0$ el error es pequeño (0.5\,\%), porque allí domina la contribución local. En
la práctica (datos de absorción medidos en una banda finita) por eso se usan modelos para las alas o
KK con sustracción.

\inc{d} \textbf{Índice de grupo y luz rápida.} Por (5.6-2), $N=n-\lambda_0\frac{dn}{d\lambda_0}$. Con
$\lambda_0=c_0/\nu$, $\frac{d}{d\lambda_0}=-\frac{\nu^2}{c_0}\frac{d}{d\nu}$, luego
$-\lambda_0\frac{dn}{d\lambda_0}=\frac{c_0}{\nu}\frac{\nu^2}{c_0}\frac{dn}{d\nu}$ y
\[
N(\nu)=n+\nu\frac{dn}{d\nu}.
\]
Estimación analítica en $\nu_0$: $n\approx1+\chi'/2$ y, derivando la forma cercana a resonancia,
$\frac{d\chi'}{d\nu}\Big|_{\nu_0}=-\chi_0\frac{\nu_0}{2}\frac{1}{(\Delta\nu/2)^2}=-\frac{2\chi_0\nu_0}{\Delta\nu^2}$, así que
\[
N(\nu_0)\approx1-\chi_0\frac{\nu_0^2}{\Delta\nu^2}=1-\chi_0Q^2=1-10^{-4}\times10^4=0 .
\]
\begin{figure}[H]\centering\includegraphics[width=0.72\linewidth]{t2_p13_grupo.png}
\caption{Problema 13(d): índice de grupo cerca de la resonancia; sombreada, la región de luz rápida.}\end{figure}
El cálculo numérico da $N(\nu_0)=0.0004$, $N<1$ entre \SI{497.55}{THz} y \SI{502.50}{THz}
($\approx\nu_0\pm\Delta\nu/2$, donde $dn/d\nu<0$: dispersión anómala) y máximos $N\approx1.13$ a los lados.
\resultado{Luz ``rápida'' ($N<1$, $v=c_0/N>c_0$) en la banda $\nu_0\pm\Delta\nu/2$; con estos parámetros
$N$ baja justo hasta $\approx0$ ($v\to\infty$), y sería negativo si $\chi_0Q^2>1$.}
\emph{¿Viola la causalidad?} No (Saleh, pp.~188--189). (1) Justo en esa banda la absorción es máxima:
$\alpha\sim\SI{5e4}{}$--$\SI{1e5}{m^{-1}}$, de modo que la luz recorre unas decenas de micras antes de
extinguirse, y el pulso se deforma: el concepto de velocidad de grupo (Taylor a primer orden) ya no
describe la propagación. (2) La velocidad de grupo es la del pico de una envolvente suave, no la de
la información; un frente abrupto (punto no analítico) nunca viaja más rápido que $c_0$. (3) De hecho
el propio $\chi(\nu)$ usado satisface Kramers--Kronig, que es la expresión matemática de la causalidad:
el modelo es causal por construcción.

\idea{Lorentz: $\chi''$ lorentziana de FWHM $\Delta\nu=\sigma/2\pi$ y pico $\chi_0Q$; $\chi'$ dispersiva con
extremos $\pm\chi_0Q/2$ en $\nu_0\mp\Delta\nu/2$. En Kramers--Kronig verifique el signo con el límite
estático ($\chi'(0)>0$). $N=n+\nu\,dn/d\nu$: luz rápida donde $dn/d\nu<0$ (dentro de la línea de absorción).}

%=====================================================================
\sol{14}{Onda modulada en amplitud en un medio dispersivo}{Saleh Prob.~5.5-2, p.~195; §22.3A, ecs.~22.3-1 a 22.3-3, pp.~961--962}\label{sol:14}

\textbf{Paso 1: descomposición espectral.} Con $\cos x=\tfrac12(e^{jx}+e^{-jx})$,
\[
\mathcal A(0,t)e^{j2\pi\nu_0t}=e^{j2\pi\nu_0t}+\frac m2e^{j2\pi(\nu_0+f_s)t}+\frac m2e^{j2\pi(\nu_0-f_s)t}:
\]
tres ondas monocromáticas (portadora y dos bandas laterales).

\textbf{Paso 2: propagación.} En un medio sin pérdidas cada componente de frecuencia $\nu$ se multiplica
por $e^{-j\beta(\nu)z}$ (función de transferencia, ec.~22.3-1, p.~961). Con $\beta(\nu_0)=\beta_0$,
$\beta(\nu_0-f_s)=\beta_1$ y $\beta(\nu_0+f_s)=\beta_2$:
\[
U(z,t)=e^{j2\pi\nu_0t}\left[e^{-j\beta_0z}+\frac m2e^{j2\pi f_st-j\beta_2z}+\frac m2e^{-j2\pi f_st-j\beta_1z}\right].
\]
\textbf{Paso 3: envolvente.} Factorizando la fase de la portadora, $U=\mathcal A(z,t)e^{-j\beta_0z}e^{j2\pi\nu_0t}$ con
\[
\mathcal A(z,t)=1+\frac m2\left[e^{j[2\pi f_st-(\beta_2-\beta_0)z]}+e^{-j[2\pi f_st+(\beta_1-\beta_0)z]}\right].
\]
Se definen
\[
\Phi=\frac{(\beta_1+\beta_2-2\beta_0)z}{2},\qquad \Psi=\frac{(\beta_2-\beta_1)z}{2},
\]
de modo que $(\beta_2-\beta_0)z=\Phi+\Psi$ y $(\beta_1-\beta_0)z=\Phi-\Psi$. Entonces
\[
\mathcal A(z,t)=1+\frac m2e^{-j\Phi}\left[e^{j(2\pi f_st-\Psi)}+e^{-j(2\pi f_st-\Psi)}\right]
=1+m\,e^{-j\Phi}\cos(2\pi f_st-\Psi).
\]
\resultado{$\mathcal A(z,t)=1+m\,e^{-j\Phi}\cos(2\pi f_st-\Psi)$, con $\Phi=(\beta_1+\beta_2-2\beta_0)z/2$ y
$\Psi=(\beta_2-\beta_1)z/2$.}

\textbf{Interpretación.} $\Psi=\frac{\beta_2-\beta_1}{2f_s}f_sz\approx\frac{d\beta}{d\nu}f_sz=2\pi f_s\frac{z}{v}$:
es el retardo de grupo de la modulación, $\cos[2\pi f_s(t-z/v)]$, no deforma nada. $\Phi$ mide la
curvatura de $\beta(\nu)$ (la GVD). Si $\Phi=q\pi$ ($q$ entero), $e^{-j\Phi}=(-1)^q$ es real:
\[
\mathcal A=1\pm m\cos(2\pi f_st-\Psi),
\]
\emph{modulación de amplitud pura, sin modulación de fase} (con la modulación invertida si $q$ es impar).
En cambio, si $\Phi=\pi/2+q\pi$, $\mathcal A=1\mp jm\cos(\ldots)$: casi pura modulación de fase, y
$|\mathcal A|^2=1+m^2\cos^2(\ldots)$ no tiene componente a $f_s$: un fotodetector no ve el tono de RF
(``desvanecimiento'' de la portadora).

\textbf{Distancias en función de $\beta''$.} Taylor en $\omega$: $\beta(\omega_0\pm\Omega)=\beta_0\pm\beta'\Omega+\tfrac12\beta''\Omega^2$,
con $\Omega=2\pi f_s$. Sumando las dos bandas laterales, los términos lineales se cancelan:
$\beta_1+\beta_2-2\beta_0=\beta''\Omega^2=4\pi^2\beta''f_s^2$. La condición $\Phi=q\pi$ es
\[
\frac{4\pi^2\beta''f_s^2\,z}{2}=q\pi
\quad\Longrightarrow\quad
z_q=\frac{2q\pi}{\beta_1+\beta_2-2\beta_0}\approx\frac{q}{2\pi\beta''f_s^2}.
\]
\resultado{AM pura en $z_q=q/(2\pi|\beta''|f_s^2)$, $q=1,2,\ldots$; desvanecimiento total del tono
detectado en los puntos intermedios $z_{q+1/2}$.}
\emph{Orden de magnitud.} Fibra estándar a \SI{1.55}{\um}, $|\beta''|=\SI{21.7}{ps^2/km}=\SI{2.17e-26}{s^2/m}$
y $f_s=\SI{10}{GHz}$: $z_1=1/[2\pi(2.17\times10^{-26})(10^{20})]=\SI{73}{km}$ y el primer desvanecimiento a
\SI{37}{km}. Por eso los enlaces de radio sobre fibra a decenas de GHz usan modulación de banda lateral
única o compensación de dispersión.

\idea{La GVD actúa sobre la \emph{suma} $\beta_1+\beta_2-2\beta_0$ (segunda diferencia), mientras que el
retardo depende de la diferencia $\beta_2-\beta_1$. La dispersión convierte AM en PM y de vuelta,
periódicamente en $z$. No confunda $d^2\beta/d\nu^2=4\pi^2\beta''$ con $\beta''=d^2\beta/d\omega^2$.}

%=====================================================================
\sol{15}{Sellmeier, índice de grupo y GVD}{Saleh Prob.~5.6-1, p.~196; ec.~5.5-28, p.~180; ecs.~5.6-2, 5.6-4, 5.6-7, pp.~185--186; Ej.~5.6-1, p.~189}\label{sol:15}

\inc{a} \textbf{Expresiones analíticas.} Sellmeier (5.5-28): $n^2=1+\sum_iB_i\frac{\lambda^2}{\lambda^2-\lambda_i^2}$
(con $\lambda\equiv\lambda_0$, $B_i\equiv\chi_{0i}$). \emph{Primera derivada:}
\[
\frac{d}{d\lambda}\frac{\lambda^2}{\lambda^2-\lambda_i^2}=\frac{2\lambda(\lambda^2-\lambda_i^2)-\lambda^2(2\lambda)}{(\lambda^2-\lambda_i^2)^2}
=-\frac{2\lambda\lambda_i^2}{(\lambda^2-\lambda_i^2)^2},
\qquad
2n\frac{dn}{d\lambda}=\frac{d(n^2)}{d\lambda}
\ \Longrightarrow\
\frac{dn}{d\lambda}=-\frac{\lambda}{n}\sum_i\frac{B_i\lambda_i^2}{(\lambda^2-\lambda_i^2)^2}.
\]
\emph{Índice de grupo} (5.6-2):
\[
\boxed{N=n-\lambda\frac{dn}{d\lambda}=n+\frac{\lambda^2}{n}\sum_i\frac{B_i\lambda_i^2}{(\lambda^2-\lambda_i^2)^2}}
\]
\emph{Segunda derivada.} Derivando $2nn'=(n^2)'$ otra vez: $2n'^2+2nn''=(n^2)''$, luego
$n''=\frac{(n^2)''-2n'^2}{2n}$, con
\[
(n^2)''=\frac{d}{d\lambda}\left[-2\lambda\sum_i\frac{B_i\lambda_i^2}{(\lambda^2-\lambda_i^2)^2}\right]
=\sum_iB_i\lambda_i^2\left[\frac{-2}{(\lambda^2-\lambda_i^2)^2}+\frac{8\lambda^2}{(\lambda^2-\lambda_i^2)^3}\right]
=\sum_i\frac{B_i\lambda_i^2(6\lambda^2+2\lambda_i^2)}{(\lambda^2-\lambda_i^2)^3}.
\]
\emph{Coeficiente de dispersión} (5.6-7):
\[
\boxed{\begin{aligned}D_\lambda&=-\frac{\lambda}{c_0}\frac{d^2n}{d\lambda^2}\\
&=-\frac{\lambda}{2nc_0}\left[\sum_i\frac{B_i\lambda_i^2(6\lambda^2+2\lambda_i^2)}{(\lambda^2-\lambda_i^2)^3}-\frac{2\lambda^2}{n^2}\Big(\sum_i\frac{B_i\lambda_i^2}{(\lambda^2-\lambda_i^2)^2}\Big)^2\right]\end{aligned}}
\]
Un término con $\lambda_i=0$ (como el primero del GaAs) solo aporta una constante a $n^2$ y no entra
en las derivadas. El programa verifica estas fórmulas contra diferencias finitas (coinciden en 6--7 cifras).

\inc{b} \textbf{[N] Sílice fundida} ($\lambda_i=0.06840$, $0.1162$, $9.8962$ \si{\um}; $B_i=0.6962$, $0.4079$, $0.8975$).
\codigo{t2_p15_sellmeier.py}
\begin{figure}[H]\centering\includegraphics[width=0.78\linewidth]{t2_p15_silice.png}
\caption{Problema 15(b): $n$, $N$ y $D_\lambda$ de la sílice fundida entre 0.25 y \SI{3.5}{\um}.}\end{figure}
\begin{center}\small
\begin{tabular}{ccccc}\toprule
$\lambda_0$ (\si{\um}) & $n$ & $N$ & $D_\lambda$ (ps/(km$\cdot$nm)) & $\beta_2$ (ps$^2$/km $=$ fs$^2$/mm)\\\midrule
0.60 & 1.45803 & 1.47801 & $-292.4$ & $+55.9$\\
0.80 & 1.45331 & 1.46713 & $-106.4$ & $+36.1$\\
1.00 & 1.45041 & 1.46303 & $-39.8$ & $+21.1$\\
1.30 & 1.44691 & 1.46162 & $+2.7$ & $-2.4$\\
1.55 & 1.44402 & 1.46259 & $+21.9$ & $-28.0$\\\bottomrule
\end{tabular}\end{center}
Se reproduce la Fig.~5.6-5 (p.~189): $n$ decrece monótonamente, $N$ tiene un mínimo ($v$ máxima) y
$D_\lambda$ pasa de unos $-200$ ps/(km$\cdot$nm) cerca de \SI{0.65}{\um} a $+20$ cerca de \SI{1.55}{\um}.
\resultado{$D_\lambda=0$ en $\lambda_{ZD}=\SI{1.273}{\um}$, donde $N$ es mínimo ($N=1.4616$). Dispersión
normal ($D_\lambda<0$) por debajo y anómala por encima.}
El libro dice \SI{1.276}{\um}; con sus mismos coeficientes (y también con los de Malitson completos)
sale \SI{1.2727}{\um}: la diferencia de 3 nm está en el texto del libro y no cambia nada físico.

\inc{c} \textbf{[N] GaAs} ($\lambda_i=0$, $0.4082$, $37.17$ \si{\um}; $B_i=3.5$, $7.4969$, $1.9347$), mismo programa.
\begin{figure}[H]\centering\includegraphics[width=0.78\linewidth]{t2_p15_gaas.png}
\caption{Problema 15(c): $n$, $N$ y $D_\lambda$ del GaAs entre 1.5 y \SI{10.5}{\um}.}\end{figure}
\begin{center}\small
\begin{tabular}{ccccc}\toprule
$\lambda_0$ (\si{\um}) & $n$ & $N$ & $D_\lambda$ (ps/(km$\cdot$nm)) & $\beta_2$ (ps$^2$/km)\\\midrule
1.5 & 3.5487 & 3.7321 & $-1324$ & $+1582$\\
2.0 & 3.5096 & 3.6081 & $-506$ & $+1074$\\
4.0 & 3.4718 & 3.5013 & $-52.3$ & $+444$\\
6.0 & 3.4612 & 3.4867 & $-7.4$ & $+142$\\
10.0 & 3.4436 & 3.4945 & $+17.3$ & $-918$\\\bottomrule
\end{tabular}\end{center}
\resultado{GaAs: $D_\lambda=0$ en $\lambda_{ZD}=\SI{6.84}{\um}$ ($D_\lambda=-506$ ps/(km$\cdot$nm) a \SI{2}{\um}, $+17.3$ a \SI{10}{\um}).}
\emph{Comparación.} (1) El GaAs tiene $n\approx3.5$ frente a $1.45$: su resonancia electrónica
(\SI{0.41}{\um}, cerca del borde de banda prohibida, \SI{0.87}{\um}) es mucho más fuerte y está más cerca del
IR. (2) Por eso, a la misma $\lambda$, su dispersión normal es enorme (a \SI{1.5}{\um} $|D_\lambda|$ es
unas 70 veces la de la sílice, que allí vale $\approx+18$, y $N-n=0.18$ frente a $0.02$). (3) El cero de dispersión cae donde se
equilibran la curvatura de la resonancia UV y la de la resonancia de red (IR): en la sílice ambas
están en $0.07$--$0.12$ y \SI{9.9}{\um}, y el cero cae en \SI{1.27}{\um}; en el GaAs, en $0.41$ y \SI{37}{\um},
y el cero se corre a \SI{6.8}{\um}. (4) Para óptica no lineal: la sílice permite dispersión anómala
(solitones, supercontinuo) a \SI{1.55}{\um}; el GaAs masivo es fuertemente normal en el IR cercano y
hay que corregirlo con la geometría de la guía.

\inc{d} \textbf{$\beta_2$ a \SI{1.55}{\um} y relación con $D_\lambda$.} De (5.6-4), $D_\nu=2\pi\beta''$, y de
$D_\lambda\dd\lambda=D_\nu\dd\nu$ con $d\nu/d\lambda=-c_0/\lambda^2$: $D_\lambda=-\frac{c_0}{\lambda^2}D_\nu$. Por tanto
\[
\beta_2=\beta''=\frac{D_\nu}{2\pi}=-\frac{\lambda^2}{2\pi c_0}D_\lambda
=-\frac{(1.55\times10^{-6})^2(21.92\times10^{-6}\ \si{s/m^2})}{2\pi(2.998\times10^8)}=\SI{-2.796e-26}{s^2/m}.
\]
\resultado{$\beta_2=\SI{-28.0}{ps^2/km}$ a \SI{1.55}{\um} (dispersión anómala). Regla:
$\beta_2\,[\si{ps^2/km}]=-1.275\,D_\lambda\,[\si{ps/(km.nm)}]$ a \SI{1.55}{\um}. La fibra estándar tiene
$D_\lambda\approx17$ (no 21.9) porque la dispersión de guía de onda resta a la del material.}
(Unidades: $1\ \si{ps/(km.nm)}=\SI{e-6}{s/m^2}$ y $1\ \si{ps^2/km}=\SI{e-27}{s^2/m}=1\ \si{fs^2/mm}$.)

\idea{$N=n-\lambda\,n'$ y $D_\lambda=-(\lambda/c_0)n''$. Signos: normal $\Leftrightarrow\beta_2>0\Leftrightarrow D_\nu>0\Leftrightarrow D_\lambda<0$.
Convierta siempre con $\beta_2=-\lambda^2D_\lambda/2\pi c_0$, y recuerde que ps$^2$/km y fs$^2$/mm son la misma unidad.}

%=====================================================================
\sol{16}{Índice del aire; velocidad de grupo en un metal}{Saleh Probs.~5.6-2, 5.6-3, p.~196; ecs.~5.6-2, 5.6-7, pp.~185--186; ec.~5.5-39, p.~183}\label{sol:16}

\inc{a} \textbf{Ajuste parabólico.} Con $x=\lambda-0.81$ \si{\um}, se busca
$n-1=a_0+a_1x+a_2x^2$ por los tres puntos $x=-0.05,0,+0.05$:
\[
a_0=2.669\times10^{-4},\quad
a_1=\frac{(2.665-2.672)\times10^{-4}}{0.10}=\SI{-7.0e-6}{\um^{-1}},
\]
\[
a_2=\frac{(2.665-2\times2.669+2.672)\times10^{-4}}{2(0.05)^2}=\SI{-2.0e-5}{\um^{-2}}
\]
(diferencias centradas, exactas para una parábola). Velocidad de grupo con (5.6-2):
\[
v(\lambda)=\frac{c_0}{N(\lambda)},\qquad
N(\lambda)=1+a_0+a_1x+a_2x^2-\lambda\,(a_1+2a_2x),\quad x=\lambda-\SI{0.81}{\um}.
\]
En \SI{0.81}{\um}: $N=1.0002669+0.81\times7.0\times10^{-6}=1.0002726$ y $v=\SI{2.997183e8}{m/s}$.
En 0.76 y \SI{0.86}{\um}: $N=1.0002710$ y $1.0002742$. Dispersión (5.6-7) con $n''=2a_2=\SI{-4.0e-5}{\um^{-2}}=\SI{-4.0e7}{m^{-2}}$:
\[
D_\lambda=-\frac{\lambda}{c_0}n''=-\frac{0.81\times10^{-6}}{2.998\times10^8}(-4.0\times10^{7})=\SI{1.08e-7}{s/m^2}=\SI{0.108}{ps/(km.nm)}.
\]
\codigo{t2_p16_aire.py}
\begin{figure}[H]\centering\includegraphics[width=0.7\linewidth]{t2_p16_aire.png}
\caption{Problema 16(a): ajuste parabólico de $n-1$ del aire e índice de grupo $N-1$.}\end{figure}
\resultado{$N(\SI{0.81}{\um})=1.000273$, $v\approx c_0(1-2.73\times10^{-4})$; $D_\lambda\approx+0.108$ ps/(km$\cdot$nm),
unas 200 veces menor que la de la fibra de sílice a \SI{1.55}{\um} (21.9). (Con solo tres puntos, $n''$ es
constante: el ajuste no puede decir cómo varía $D_\lambda$; la variación que imprime el programa, de
$0.101$ a $0.115$, viene solo del factor $\lambda$.)}
\emph{Comentario sobre el signo.} $n''<0$ aquí, mientras que la sílice en el visible tiene $n''>0$: con tres
datos de 4 cifras, la segunda diferencia es $-0.001\times10^{-4}$, del orden del redondeo del último dígito,
así que el signo de $D_\lambda$ no es confiable; lo que sí es robusto es el orden de magnitud: el aire es
prácticamente no dispersivo ($|D_\lambda|\lesssim0.1$ ps/(km$\cdot$nm)).

\inc{b} \textbf{Metal de Drude: $v_fv_g=c_0^2$.} De (5.5-39), $\epsilon_e=\epsilon_0(1-\omega_p^2/\omega^2)$. Para
$\omega>\omega_p$:
\[
\beta(\omega)=\omega\sqrt{\epsilon_e\mu_0}=\frac{\omega}{c_0}\sqrt{1-\frac{\omega_p^2}{\omega^2}}=\frac{\sqrt{\omega^2-\omega_p^2}}{c_0}.
\]
Velocidad de fase: $v_f=\omega/\beta=\dfrac{c_0\,\omega}{\sqrt{\omega^2-\omega_p^2}}$ ($>c_0$). Velocidad de grupo (5.6-1):
\[
\frac{1}{v_g}=\frac{d\beta}{d\omega}=\frac{1}{c_0}\frac{\omega}{\sqrt{\omega^2-\omega_p^2}}
\quad\Longrightarrow\quad v_g=\frac{c_0\sqrt{\omega^2-\omega_p^2}}{\omega}\ (<c_0).
\]
\resultado{$v_f\,v_g=\dfrac{c_0\omega}{\sqrt{\omega^2-\omega_p^2}}\cdot\dfrac{c_0\sqrt{\omega^2-\omega_p^2}}{\omega}=c_0^2$.}
Equivalente: $\omega^2=\omega_p^2+c_0^2\beta^2$ (relación de dispersión hiperbólica, Fig.~5.5-10, p.~183) da
$2\omega\dd\omega=2c_0^2\beta\dd\beta$, o sea $(\omega/\beta)(d\omega/d\beta)=c_0^2$. Así la fase va más rápido que
$c_0$ pero la energía e información viajan a $v_g<c_0$.

\idea{Para $N$ use la primera derivada y para $D_\lambda$ la segunda; con datos discretos, las diferencias
centradas son exactas para una parábola. En un plasma $n<1$ pero $N=1/n>1$.}

%=====================================================================
\sol{17}{Pulsos gaussianos: ensanchamiento, chirp y compresión}{Saleh §5.6, pp.~184--187; §22.1, ecs.~22.1-4, 22.1-5, 22.1-8, Tabla 22.1-1, pp.~939--943; §22.3, ecs.~22.3-3 a 22.3-18, Tabla 22.3-1, pp.~961--966}\label{sol:17}

\textbf{Fórmulas base.} Un pulso gaussiano sin chirp $\mathcal A(0,t)=A_0e^{-t^2/\tau_0^2}$ en un medio con
$D_\nu=2\pi\beta''$ sigue siendo gaussiano con (Tabla 22.3-1, p.~964)
\[
\tau(z)=\tau_0\sqrt{1+(z/z_0)^2},\qquad a(z)=\frac{z}{z_0},\qquad z_0=\frac{\pi\tau_0^2}{D_\nu}=\frac{\tau_0^2}{2\beta''},
\]
y el centro llega con retardo $z/v=Nz/c_0$. Como $\tau_{\rm FWHM}\propto\tau$, la misma ley vale para el FWHM.

\inc{a} \textbf{100 fs y 20 fs en 1 cm de sílice a 800 nm.} Del programa del Prob.~15:
$n=1.45331$, $N=1.46713$. Retardo de grupo
\[
\tau_d=\frac{NL}{c_0}=\frac{1.46713\times\SI{0.01}{m}}{\SI{2.998e8}{m/s}}=\SI{48.94}{ps}.
\]
GVD: $\beta''=-\frac{\lambda^2}{2\pi c_0}D_\lambda$ con $D_\lambda=-106.4$ ps/(km$\cdot$nm) da
$\beta''=\SI{36.14}{fs^2/mm}$, y $D_\nu=2\pi\beta''=\SI{227.0}{fs^2/mm}=\SI{2.27e-25}{s^2/m}$.
\begin{itemize}
\item $\tau_{\rm FWHM}=\SI{100}{fs}$: $\tau_0=100/1.1774=\SI{84.93}{fs}$, $\tau_0^2=\SI{7213}{fs^2}$,
$z_0=\tau_0^2/2\beta''=7213/72.28=\SI{99.8}{mm}$. Tras $L=\SI{10}{mm}$: $L/z_0=0.100$ y
$\tau_{\rm FWHM}=100\sqrt{1+0.0100}=\SI{100.5}{fs}$.
\item $\tau_{\rm FWHM}=\SI{20}{fs}$: $\tau_0=\SI{16.99}{fs}$, $\tau_0^2=\SI{288.5}{fs^2}$, $z_0=\SI{3.99}{mm}$; $L/z_0=2.505$,
$\tau_{\rm FWHM}=20\sqrt{1+6.275}=\SI{53.9}{fs}$, con chirp $a=2.5$.
\end{itemize}
\resultado{$N=1.4671$; retardo \SI{48.9}{ps}; $\beta''=\SI{36.1}{fs^2/mm}$; 100 fs: $z_0=\SI{99.8}{mm}$, sale con
\SI{100.5}{fs}; 20 fs: $z_0=\SI{3.99}{mm}$, sale con \SI{53.9}{fs}.}
\emph{¿Por qué tan distinto?} $z_0\propto\tau_0^2$: reducir el pulso 5 veces reduce la longitud de
dispersión 25 veces. Físicamente, el ancho espectral es $\Delta\nu=0.44/\tau_{\rm FWHM}$ (22.1-10): el pulso de
20 fs tiene \SI{22}{THz} de ancho (45 nm), y la diferencia de retardo de grupo entre sus extremos,
$D_\nu L\Delta\nu\approx\SI{50}{fs}$, supera su duración; el de 100 fs tiene 5 veces menos ancho de banda y
la diferencia de retardos (\SI{10}{fs}) es pequeña frente a 100 fs.

\inc{b} \textbf{10 ps en fibra estándar.} $D_\nu=-\frac{\lambda^2}{c_0}D_\lambda=-\frac{(1.55\times10^{-6})^2}{2.998\times10^8}(17\times10^{-6})=\SI{-1.362e-25}{s^2/m}$
($\beta_2=\SI{-21.7}{ps^2/km}$). $\tau_0=10/1.1774=\SI{8.493}{ps}$,
\[
|z_0|=\frac{\pi\tau_0^2}{|D_\nu|}=\frac{\pi(8.493\times10^{-12})^2}{1.362\times10^{-25}}=\SI{1664}{m}.
\]
$\tau(\SI{10}{km})=10\sqrt{1+(10/1.664)^2}=10\times6.094=\SI{60.9}{ps}$;
$\tau(\SI{100}{km})=10\sqrt{1+(100/1.664)^2}=\SI{601}{ps}$ (régimen lineal, $\tau\approx\tau_{\rm FWHM}z/|z_0|$).
\resultado{$D_\nu=\SI{-1.36e-25}{s^2/m}$, $|z_0|=\SI{1.66}{km}$; \SI{60.9}{ps} tras 10 km y \SI{601}{ps} tras 100 km.}

\inc{c} \textbf{¿Dónde quedan las frecuencias altas?} Tras propagarse, el chirp es $a(z)=z/z_0$ y la
frecuencia instantánea (22.1-4), (22.1-23) es
\[
\nu_i(t)=\nu_0+\frac{1}{2\pi}\frac{d\varphi}{dt}=\nu_0+\frac{a(z)}{\pi\tau^2(z)}\,(t-z/v).
\]
\emph{Dispersión normal} ($D_\nu>0\Rightarrow z_0>0\Rightarrow a>0$): $\nu_i$ crece con $t$ (up-chirp): las
frecuencias bajas (rojo) quedan adelante y las altas (azul) atrás. Razón física: con $D_\nu=\frac{d}{d\nu}(1/v)>0$, el
retardo por unidad de longitud $1/v$ crece con $\nu$, así que el azul viaja más despacio y llega después.
\emph{Dispersión anómala} ($D_\nu<0$): $a<0$, down-chirp, el azul va adelante. La figura lo confirma con
el pulso de 20 fs del inciso (a) (el de anómala se calculó con $D_\nu\to-D_\nu$): las intensidades son
idénticas (el ensanchamiento no depende del signo), pero la pendiente de $\nu_i(t)$ se invierte.
\begin{figure}[H]\centering\includegraphics[width=0.78\linewidth]{t2_p17_chirp.png}
\caption{Problema 17(c): pulso de 20 fs tras 1 cm con $|D_\nu|$ de la sílice a 800 nm. Arriba, intensidad;
abajo, frecuencia instantánea $\nu_i-\nu_0$ (solo donde la intensidad supera el 2\,\% del pico).}\end{figure}
\resultado{Normal: altas frecuencias en la cola (up-chirp, $a>0$). Anómala: altas frecuencias al frente
(down-chirp, $a<0$). Ojo: en la convención $e^{-i\omega t}$ (Boyd, Agrawal) el signo de la fase del chirp es
el opuesto para el mismo pulso físico.}

\inc{d} \textbf{[N] Propagación por FFT.} Se implementa $A(z,t)=\mathcal F^{-1}\{e^{-j\pi D_\nu f^2z}\mathcal F\{A(0,t)\}\}$
(ec.~22.3-3 con $b=D_\nu z/\pi$, en el marco que viaja a $v$), con \texttt{numpy.fft}, cuyo signo
coincide con el de Saleh.
\codigo{t2_p17_pulsos.py}
\emph{Verificación de (a):} la FFT da \SI{100.50}{fs} y \SI{53.94}{fs}, igual a la fórmula analítica en 4 cifras.

\emph{Pulso con chirp inicial.} En unidades con $\tau_0=1$ y $z_0'\equiv\pi\tau_0^2/|D_\nu|=1$ (longitud de
dispersión con el ancho de \emph{entrada}), se toma el pulso del enunciado
$A(0,t)=\exp[-(1+ja)t^2/\tau_0^2]$ con $a=-3$, es decir
\[
A(0,t)=e^{-t^2/\tau_0^2}\,e^{+j3t^2/\tau_0^2},\qquad \varphi(t)=+\frac{3t^2}{\tau_0^2},\qquad
\nu_i=\nu_0+\frac{3t}{\pi\tau_0^2}.
\]
\textbf{Atención:} en la convención de Saleh ($e^{j2\pi\nu_0t}$) este pulso tiene $\nu_i$ \emph{creciente}:
es up-chirp, con parámetro de Saleh $a_1=+3$ en $\exp[-(1-ja_1)t^2/\tau^2]$ (22.1-12). Por la regla de
Saleh (p.~967), un pulso up-chirp se comprime solo en un medio \emph{anómalo}; en uno normal se ensancha
desde el principio. Analíticamente, con $\mathcal A\propto e^{-\gamma t^2}$, $\gamma=(1-ja_1)/\tau_0^2$, su
transformada es $\propto e^{-\pi^2f^2/\gamma}$; el filtro suma $j D_\nu z/\pi$ a $1/\gamma$:
\[
\frac{1}{\gamma(z)}=\frac{\tau_0^2(1+ja_1)}{1+a_1^2}+j\frac{D_\nu z}{\pi}.
\]
El pulso es mínimo (sin chirp, $1/\gamma$ real) cuando la parte imaginaria se anula:
\[
z_{\min}=-\frac{a_1}{1+a_1^2}\frac{\pi\tau_0^2}{D_\nu}=-\frac{a_1}{1+a_1^2}z_0',
\qquad
\tau_{\min}=\sqrt{\frac{\tau_0^2}{1+a_1^2}}=\frac{\tau_0}{\sqrt{1+a_1^2}},
\]
que es (22.3-14)--(22.3-16) expresado con el ancho de entrada ($z_0$ del libro usa el ancho mínimo:
$z_0=z_0'/(1+a_1^2)$). Con $a_1=+3$, $z_{\min}>0$ exige $D_\nu<0$.
\begin{figure}[H]\centering\includegraphics[width=0.78\linewidth]{t2_p17_compresion.png}
\caption{Problema 17(d): ancho del pulso (FWHM normalizado a $1.177\tau_0$) contra la distancia, para el
pulso del enunciado en dispersión normal y anómala y para el pulso con el chirp opuesto en dispersión
normal.}\end{figure}
Resultados de la FFT (comparados con $\tau_0/\sqrt{10}=0.3162\tau_0$ y $3/10=0.3$):
\begin{center}\small
\begin{tabular}{lcc}\toprule
Caso & ancho mínimo & distancia \\\midrule
Enunciado literal ($a=-3$, up-chirp), dispersión normal & $1.000\,\tau_0$ (en $z=0$) & no se comprime\\
Enunciado literal ($a=-3$, up-chirp), dispersión anómala & $0.3163\,\tau_0$ & $0.300\,z_0'$\\
$a=+3$ en la forma del enunciado (down-chirp), normal & $0.3163\,\tau_0$ & $0.300\,z_0'$\\\bottomrule
\end{tabular}\end{center}
\resultado{Un pulso down-chirp en dispersión normal (o up-chirp en anómala) con $|a|=3$ se comprime a
$\tau_{\min}=\tau_0/\sqrt{1+a^2}=0.316\tau_0$ a la distancia $z=\frac{|a|}{1+a^2}z_0'=0.3\,z_0'$, con
$z_0'=\pi\tau_0^2/|D_\nu|$ calculado con el ancho de entrada. Tal como está escrito el enunciado
($a=-3$ en $e^{-(1+ja)t^2/\tau_0^2}$, con la propagación de Saleh), el pulso es up-chirp y en dispersión
normal \emph{no} se comprime; se comprime igual (0.316$\tau_0$ a $0.3z_0'$) en dispersión anómala.}
\emph{Sobre la discrepancia con la respuesta publicada.} La respuesta dice ``$a=3$ es down-chirp en
$e^{j\omega t}$; se comprime en dispersión normal''. Eso es correcto para $a=+3$ en la forma
$e^{-(1+ja)t^2/\tau_0^2}$ (fase $-3t^2/\tau_0^2$), pero el enunciado pone $a=-3$. La forma $(1+ja)$ con $a<0$
como pulso ``down-chirp'' es la de Agrawal, que usa $e^{-i\omega_0t}$; al combinarla con el propagador de Saleh
el signo se invierte. Los valores numéricos ($0.3z_0'$, $0.316\tau_0$) son los mismos en ambos casos.

\idea{Tres reglas: (1) $z_0=\pi\tau_0^2/|D_\nu|$ con $\tau_0$ de $1/e$ en amplitud (FWHM$/1.177$), no el FWHM.
(2) Normal $\Rightarrow$ up-chirp $\Rightarrow$ el azul atrás. (3) Hay compresión si el chirp del pulso y
el que impone el medio tienen signo opuesto; el ancho mínimo es el limitado por transformada,
$\tau_0/\sqrt{1+a^2}$. Antes de decidir si un pulso es up o down, calcule $\nu_i=\nu_0\pm\frac{1}{2\pi}\frac{d\varphi}{dt}$
con el signo de la convención del texto ($+$ en Saleh, $-$ en Boyd).}
